\chapter*{Glosarium}
\addcontentsline{toc}{chapter}{Glosarium}

Istilah asing dicetak miring di dalam teks. Daftar berikut disusun menurut abjad.

\begin{description}[leftmargin=4.4cm, style=nextline, itemsep=6pt]
  \item[Halusinasi (\textit{hallucination})] Keluaran model yang terdengar meyakinkan tetapi tidak sesuai dengan isi rekaman, misalnya kalimat yang berulang tanpa dasar.
  \item[Kecerdasan artifisial (\textit{artificial intelligence}, AI)] Teknologi yang memungkinkan komputer menyusun, merangkum, atau mengenali pola dari data. Pandora memakainya untuk merapikan teks dan menyusun naskah.
  \item[Kepala dokumen] Bagian awal naskah dinas yang memuat tanggal, tempat, unit, dan keterangan lain yang diisi pengguna. Pandora tidak mengarang isi kolom ini.
  \item[Keluaran] Berkas hasil yang dibuat Pandora, misalnya transkrip, notulensi, atau naskah dalam format tertentu.
  \item[Model] Program pengenal suara yang mengubah audio menjadi teks. Makin besar model, makin teliti, tetapi prosesnya makin lama.
  \item[Naskah dinas] Dokumen resmi instansi dengan bagian dan bahasa baku, misalnya laporan kegiatan atau notulen rapat.
  \item[Notulensi (\textit{minutes})] Catatan resmi jalannya rapat yang memuat ringkasan, pembahasan, keputusan, dan tindak lanjut.
  \item[Rapikan transkrip] Fungsi untuk menyusun ulang transkrip agar lebih mudah dibaca. Transkrip asli tidak berubah.
  \item[Riwayat] Daftar folder hasil milik akun Anda yang dapat dibuka kembali.
  \item[Segmen] Potongan teks pendek dengan waktu awal dan akhir dalam rekaman.
  \item[Transkrip] Teks hasil pengubahan rekaman audio.
  \item[Transkripsi (\textit{transcription})] Proses mengubah rekaman suara menjadi teks.
  \item[Unggah] Tindakan memasukkan berkas audio dari komputer ke Pandora.
\end{description}
